\section{Methodology}
\label{S:methodology}

% --- Replace everything below with your own content. ---

\subsection{Dataset}\label{SS:dataset}


We use a dataset consisting of \textit{[N observations/records]} collected from \textit{[source, e.g., a public repository, a survey, a lab experiment]}. Each observation includes the following variables:

\begin{itemize}
    \item \textbf{[Variable 1]} --- \textit{[type: continuous / categorical
    / binary; brief description]}
    \item \textbf{[Variable 2]} --- \textit{[type; description]}
    \item \textbf{[Outcome variable]} --- \textit{[type; description]}
\end{itemize}

Raw data is stored under \texttt{data/}; any cleaning or preprocessing steps (e.g., removing missing values, outlier filtering, normalization) should be scripted under \texttt{code/} so the pipeline is reproducible rather than performed manually.



\subsection{Statistical Methods}

To test our hypotheses, we use \textit{[e.g., a two-sample $t$-test / one-way ANOVA / linear regression /chi-squared test]}. We report effect sizes (e.g., Cohen's $d$) alongside $p$-values, and we set the significance threshold at $\alpha = 0.05$ unless stated otherwise. Where multiple comparisons are performed, we apply \textit{[e.g., a Bonferroni correction]} to control the family-wise error rate.

\begin{figure}[H]
    \centering
    
% figures/example-tikz.tex
%
% A standalone TikZ diagram, kept separate from the section files so that
% the .tex sources stay readable. Include it inside a figure environment,
% e.g. from within a section file:
%
%   \begin{figure}[H]
%       \centering
%       \input{figures/example-tikz}
%       \caption{Randomized two-group experimental design.}
%       \label{fig:rct-design}
%   \end{figure}
%
% Requires (already loaded in main.tex):
%   \usepackage{tikz}
%   \usetikzlibrary{positioning, arrows.meta, shapes.geometric}

\begin{tikzpicture}[
    node distance = 1.2cm and 1.8cm,
    box/.style = {
        draw, rounded corners, minimum width=3cm, minimum height=1cm,
        align=center, font=\small
    },
    decision/.style = {
        draw, diamond, aspect=2, minimum width=2.6cm, minimum height=1cm,
        align=center, font=\small, inner sep=1pt
    },
    arrow/.style = {-{Latex[length=2mm]}, thick}
]

    % --- Nodes ---
    \node[box] (population) {Target Population\\($N$ units)};
    \node[decision, below=of population] (randomize) {Random\\Assignment};

    \node[box, below left=1.4cm and 1.2cm of randomize] (control)
        {Control Group\\(no treatment)};
    \node[box, below right=1.4cm and 1.2cm of randomize] (treatment)
        {Treatment Group\\(receives intervention)};

    \node[box, below=1.4cm of control] (measure-c) {Measure Outcome $Y_C$};
    \node[box, below=1.4cm of treatment] (measure-t) {Measure Outcome $Y_T$};

    \node[box, below=1.6cm of randomize, yshift=-3.0cm,
          minimum width=4.2cm] (compare)
        {Compare $Y_T$ vs.\ $Y_C$\\(e.g., two-sample $t$-test)};

    % --- Arrows ---
    \draw[arrow] (population) -- (randomize);
    \draw[arrow] (randomize) -| (control);
    \draw[arrow] (randomize) -| (treatment);
    \draw[arrow] (control) -- (measure-c);
    \draw[arrow] (treatment) -- (measure-t);
    \draw[arrow] (measure-c) |- (compare);
    \draw[arrow] (measure-t) |- (compare);

\end{tikzpicture}
    \caption{Randomized two-group experimental design.}
    \label{F:rct-design}
\end{figure}

All analysis code is written in \textit{[Python]} and is available in \texttt{code/}, along with instructions for reproducing the results in \Cref{S:results}.